\section{性能分析方法论}

本章将对上一节最后一个应用程序进行一些性能分析。需要做的只是改变块大小（Block Size），看看它如何影响不同的性能指标。从 CUDA 代码开始，如所示，这里有一个将两个向量相加并将结果存储在第三个向量中的核函数，该向量称为 C 向量。然后在主函数中，遵循一般 CUDA 应用程序的步骤。如所示，定义了一些指针，然后在 GPU 中为这些向量或指针分配空间，然后在 CPU 端做同样的事情，并初始化向量 A 和 B，然后将数据从 CPU 中的 A 和 B 复制到 GPU 中的 A 和 B，以便能够在 GPU 中执行加法操作。然后执行核函数来执行加法操作，最后使用 \texttt{cudaMemcpy} API 返回结果。这就是在这个应用程序中需要做的和需要了解的全部内容。

这里需要控制或操作的是每个网格的块大小变量。实际上有两个变量，不是一个。因为如果还记得，对于任何 CUDA 核函数，有两个配置。第一个是每个块的线程数，第二个是应用程序中的块数。如果搜索这个，可以往上翻，看到正在从终端读取这个值。当执行这个应用程序时，需要在执行时添加块大小，并且这个值将被输入到这里的这条指令中。相比于上一节讨论的内容，这只是在应用程序中所做的唯一更改。

回到这里，这个文件叫 \texttt{test\_01}。这里有一个名为 Section 05 的新章节，输入 \texttt{ls} 命令，这是文件。需要编译这个文件，写 \texttt{nvcc -o}，然后需要为可执行输出起个名字。可以说例如 \texttt{test\_01}，可以就这样写，或者也可以写 \texttt{.exe}，两者都可以。然后需要 CUDA 应用程序的 \texttt{.cu} 文件，接着按回车键。这里有可执行文件，可以直接复制这个并按回车键。

需要做的就是来到这里写一个空格，然后假设块大小是 256 个线程。如所示，它没有运行，因为先去掉这个空格，还是不行。因为它在这里要求块大小。然后按回车，应用程序成功执行，没有任何问题。终端上不会显示任何内容，因为没有打印任何东西。总之这就是需要做的。可以在这里放一个 \texttt{printf} 或类似的东西，以确保执行成功完成。

\subsection{Python自动化分析脚本}

需要改变块大小，并用不同的块大小来分析这个应用程序，以观察这个变化的影响。这里有一个 Python 文件，它可以完成所有这些分析。将读取一些性能指标，然后有 L1 命中率、L2 命中率、占用率（Occupancy）。第一个是执行时间，第二个是活动周期数（Active Cycles）。因为通常可以用秒或周期数来衡量时间。这是用来读取占用率的指标。这里测量一个叫做分支效率（Branch Efficiency）的东西，将在关于分支发散（Branch Divergence）的章节中讨论它。今天不需要关注那个，但只是想提前开始了解。这个叫做分支效率。

最后有 DRAM 吞吐量。DRAM 吞吐量是通过内存每秒访问的千兆字节数（GB/s）来衡量的。需要知道这个应用程序的执行时间是多少秒，并使用该值来计算 DRAM 吞吐量。时间会影响 DRAM 吞吐量，因为吞吐量是通过每秒千兆字节数来衡量的。工具会帮做这件事，但只是在解释所有东西，这些就是指标。把它们写了两次，一次在这里，另一次在这里，但每个都有其用途，这个将和 NCU 一起使用。如果还记得，可以来到这里，像这样写 \texttt{ncu}，然后两个短横线 \texttt{--metrics}，就像在上一节讨论的那样。

然后可以在这里列出所有这些指标，如所示，得到了所有值。但在这里只是使用 Python 文件来为不同的块大小做这个，而不仅仅是一个。复制所有这些，把它们夹在这里，然后点击回车。不想看着终端显示很多表格，然后把这些值复制到 Excel 表格中再进行比较。需要 Python 脚本来完成所有这些。这就是为什么这里准备了 Python 脚本。可以用这个，但需要它们作为一个数组。如所示，使用这里的 header 变量来为从这个 Python 脚本输出的 CSV 文件创建一个表头。然后这里打开输出文件，以便在运行 NCU 时存储这些指标名称及其值。

然后对于每个块大小，从 32 开始，因为最小的块大小是一个 Warp。从 32 个线程开始，一直到 1024。一个 Warp 由 32 个线程组成，并且取最大值 1024，因为根据白皮书，不能使用超过这个数字的块大小。

如果还记得在硬件章节，是第二节或第一节。说过根据白皮书，不能使用超过 1024 个线程的块大小。然后这里创建命令，这里是 1024。因为范围在这里，所以只需要在 1024 上加 1，以确保在 for 循环中包含 1024 这个值。

这是命令，如果看这里，它和这里的命令完全一样，将使用这个变量，而不是写出所有指标。但没有尝试这些指标，因为一开始已经在这里写了，如所示。如所示，在这里使用这个变量，然后这是可执行文件。这里添加块大小，块大小如所示，范围从 32 到 1024，步长为 32。即是 32、64、96、128，以此类推，每一步增加一个 Warp 或增加 32 个线程。

然后这里这是一条 Python 指令，用于运行这个命令。这完全等同于执行这里的这个即按回车执行这个命令。这就是在 Python 编程中执行终端命令的方式。这只是一个 \texttt{print}，然后这只是一个 \texttt{print}，用于显示当前正在为这个块大小运行的信息。这意味着这里的输出或这里的表格将被存储在 result 变量中。然后当执行这个 NCU 命令时，获取输出并将其存储在 result 变量中。现在需要做的就是从这个输出中读取指标名称及其值，并将它们全部存储在一个文件中。

需要分割这个 result 变量。把它们全部存储在这里的输出文件中，在这个 CSV 文件中分成不同的行。然后遍历拥有的所有行，逐行读取。对于每一行将查看或检查指标名称。一旦有了一个指标名称，就尝试读取它的值。

一旦有了值，需要将这个值存储到输出文件中。这里有很多细节，这些细节只是为了去除所有不需要的东西。因为如果看一下输出，例如不需要所有这些。只需要指标名称和它的值。如果这样做，会发现它比这个表格更奇怪。

顺便说一下，如果重复运行，因为在这里使用两个短横线 \texttt{--csv} 来确保它们以 CSV 文件格式显示，而不是表格。看到它是什么样子了吗？需要去掉这个逗号，需要去掉这个双引号等。所有这些只是为了去除不需要的任何东西，或者只是从输出中移除它。

从这个文件中需要的就是这两样东西：指标名称和它的值。然后需要的是获取指标值，并将其存储在输出文件中。一旦在输出文件中有了所有内容，只是提一下具体存储在哪里。

是的，打开输出文件，从这里开始写表头，需要看到写入表头。看到表头了吗？因为表头应包含块大小以及每个指标的名称。是的，这里尝试获取指标值，然后将其写入输出文件。然后应该看到另一个写入指令来写入值，然后这里需要绘制每个指标随块大小变化的关系图，意思是将块大小从 32 改到 1024，与例如周期数的关系等。需要绘制所有这些，这就是在这里做的。

需要在图表中显示块大小与占用率的关系，与时间的关系，与执行时间的关系，块大小增加。

\subsection{分析结果图表}

运行，这就是全部的内容。在纵轴上绘制所有指标，在 X 轴上应该是块大小。想象一下从 32 到 1024 步长为 32，这大概有 15 次左右。

将添加这个 Python 脚本，以便不仅可以将它用于这个应用程序。因为只是给出一个如何分析许多配置的想法，只需一次运行。因为想象一下在这里运行了多少次。

来到这里使用这个命令 15 次，然后手动读取这些值并把它们放进文件，这是非常困难的，将会花费太多时间。而且如果想用不同的配置再次重复它也会非常累人。这就是 Python 文件的样子，以及这里需要做的所有事情。

需要做的就是来到这里输入 \texttt{python3}，文件名是 \texttt{get\_insight.py}，然后按回车。

它将为所有步骤或所有迭代重复运行。现在正在为块大小 32 运行。可以在这里看到这个 \texttt{print} 通过这条指令体现出来。这就是从这个脚本中需要的全部。一旦它完成了所有迭代，脚本将给出一些图表。除此之外，想再看另一个文件，需要查看图表，看看块大小的变化如何影响不同的性能指标。现在完成了，输入 \texttt{ls}，但不想让你困惑，所以等等看从这个输出中能得到什么。看看这些图表里的内容。

如所示，有一些图表，所以想到这里，在第一个图表里有 GPU 时间、执行时间和周期数。看看第一步或第一次迭代，其中块大小是 32 个线程或一个 Warp。如所示，对于两者——时间和活动周期——与其他所有迭代相比，都有更大的值。当每个块使用一个 Warp 时，得到了更长的执行时间。如在此所示，其他所有迭代中，这里只是稍微增加了一点，但大多数几乎相等。需要拿这个值和其他值进行比较，和其中任何一个值比较。这意味着需要进行一些性能分析，以理解为什么这里的执行时间比其他迭代长。

看看这如何影响其他指标。这里有 L1 缓存和 L2 缓存。对于 L1 看到对于所有迭代它都是零。因为改变块大小并不会改变访问全局内存的方式，将访问相同的扇区（Sector）或相同的字节。这就是为什么 L1 缓存是恒定的且等于 0，但在这个应用程序中没有那样做。如果改变例如步长（Stride），即每两个元素之间的步长，这可能会改变。

好的，这里 L2 大约等于 33 点几，它略有变化或变化很小，但几乎相同。在这个图表里没有看到任何变化。再次只是想说明，这里有块大小，这里是指标值，L1 命中率和 L2 命中率。因为如果看核函数时间，可以看到第一个值很大，而所有其他值都很小。

看最后一个，如所示，内存带宽与核函数时间成反比关系。对于内存带宽和 Warp 活动度或 Warp 占用率，可以在这里看到相反的情况。可以看到第一个值很小，但所有其他值都很高。那么开始做吧，这就是为什么需要关注第一个值，并拿例如第二个值对两者进行一些性能分析。来到这里说，例如需要输入 \texttt{ncu}，需要的是 \texttt{./test.exe}，然后输入 32。

需要看看当块大小是 32 个线程时，不同的部分（Section）是什么样的。网格大小（Grid Size）指的是这个应用程序中所有的块。这些是部分，需要看的是块大小，如所示是 32，然后是网格大小。这是一个很大的数字，因为有数百万个元素需要相加。如果再看这里的应用程序，可以看到这些向量中每一个的大小是这个值，这是一个非常大的值。以及块大小，如所说是 32 个线程。如果有这个数量的元素，这是一个很大的数字，而只有 32 个线程。如果用这个除以这个就得到块数，这就是得到块数的方法。

总之对于每一个网格的块数，通常使用这个公式：总大小加每个块的线程数减一，除以每个块的线程数。而这个应用程序结合给每个块的线程数的值，两者共同给出了在这个应用程序中需要的线程总数。这给了每个网格正确的块数。因为记住，想确保应用程序中的线程总数等于要相加的元素总数。

总之，这里需要做的是取这个值和这个值，并把它们放在这里。如所示，每当在应用程序中使用块大小 32 时，总块数将是这里的这个值。那么在应用程序中有多少个活动块（Active Blocks）呢？

如果还记得在上一节解释了"已分配活动块"（Allocated Active Blocks）这个术语，并且说过这真的非常重要。在继续本章之前，需要先看那个章节。总之具体在哪里可以找到这个值？有 16、128、48，可以来到这里进入占用率部分，然后取这四个值中最小的那个。已分配活动块的数量将是 16。这就是为什么这里加了 16。很好，使用 CUDA API，可以获取 GPU 上可用的 SM 数量。在 GPU 上，得到的数量是 38 个 SM。

\subsection{波次与执行时间分析}

现在想计算波次（Wave）数量，这里需要了解更多关于波次定义的信息。说过已分配活动块的数量通常指的是可以在一个 SM 中同时运行的块的数量。意思是在那种情况下，可以同时运行 16 个块在每个 SM 上。但是记住块的总数不是 16、32 甚至 64，而是超过 100 万。从那个巨大的数字中，可以在每个 SM 上同时运行 16 个，那么剩下的块怎么办？这就是所说的波次。

例如除以这个值，将这个值也就是总块数除以 16。说过可以同时运行 16 个。然后除以 SM 的数量，因为讨论的是一个 SM。这就是在应用程序中，每个 SM 将要运行的波次数量。

再重复一遍，当这个应用程序启动时，它从这个数字中取出 16 个块交给第一个 SM，意思是对于这里的每个时间点，每个波次，将在每个 SM 上运行 16 个块。同样的情况也适用于所有其他 SM。一旦这些 SM 中的任何一个完成了它的 16 个块，就意味着它完成了第一个波次。那么它会请求另一个波次，它会请求另外 16 个块。

情况就是这样，看一下这里，这是 SM，这是第一个波次。如所示，有 16 个块 ID，从 0 到 15。一旦完成了所有这些块，SM 就可以获取另一个波次，可以获取另外 16 个块。但是再次当执行第一个波次时，不能获取这些块中的任何一个。即使所有这些块都因依赖而停顿，因内存操作而停顿或因任何其他原因停顿，同时只能有 16 个块在每个 SM 上。一旦完成那个波次，可以再获取 16 个块，就是这样。

这个术语非常重要，将解释它是如何重要的。总之这是使用块大小 32 的情况。如果使用块大小 64，这些是数字，又运行了 NCU 来获取总块数，这是总块数。为什么这里的总块数是那个数字的一半？这很简单，如果增加每个块的线程数，就必须减少总块数。因为说过总块数结合每个块的大小，每个块的大小应该等于应用程序中元素的总数。

这就是它的工作原理。现在需要知道有多少个波次。再次只是想展示第一种情况的波次数量，看这里，看到这个数字了吗？每个 SM 的波次数量，可以从这里检查它，这完全是同一个数字。

回到这里除以这个值，这是总块数。应用程序中的所有块，所有这些块都应该在 38 个 SM 上执行。需要知道每个 SM 将执行多少个块，这就是需要做的。复制这个值并除以 38。那样的话，每个 SM 应该执行超过 13000 个块。

但不能同时执行所有这些块，可以在不同的波次中调度所有这些块。每个波次由 16 个块组成，需要用 16 除这个数以得到每个 SM 的波次数量。如所示是 862，如果在这里运行 NCU，会发现完全相同的数字。如所示是 862。

只是想在这里提一下波次数量，因为需要很好地理解它。总之回到第二张幻灯片，因为它非常重要。说过这种情况的执行时间比第二种情况的执行时间长。还记得在这里吗？不在这里，是在这个图表里。回头看看，还记得执行时间或活动周期数吗？对于块大小 32，它比块大小 64 的大，这没关系。

那么需要理解为什么情况是这样的。

\subsection{块大小对性能的影响}

想象有一个 GPU，两者都在执行同一个应用程序，但使用不同的块大小。想象将只分析 38 个 SM 中的一个 SM。在这种情况下有一个 SM，在这种情况下也有这个 SM，但区别在于——并且在两种情况下都有 16 个已分配活动块。确认一下。是的，如果看这里的占用率部分，可以看到这四个值的最小值是 16。在两种情况下已分配活动块的数量都是 16。这两条线代表一个 Warp，这就是为什么这里有 16 个从 0 到 15，这里也是从 0 到 15。

这里有两个 Warp，为什么？因为这个块由 64 个线程组成，块大小是 64，而一个 Warp 等于 32 个线程，这就是为什么这个块比这个大。如果用 32 除这个，那么每个块就有两个 Warp。

想象有一个内存操作（红色的那个）以及一个依赖于该内存操作的加法操作。在这种情况下，每个块将执行内存操作，并会等待内存操作完成。而内存操作通常比 GPU 中的算术操作花费更长时间。如果在 L1 缓存命中它可能花费 30 个周期，或者如果在 L2 缓存命中大约 200 个周期，或者如果在全局内存命中 300 或 400 个周期。但这里有依赖关系，一旦因这个依赖而停顿，SM 将不会执行任何操作。

所有 Warp 或所有块都会停滞。但这里有不同的 Warp。如果这里的第一个 Warp 因任何原因停顿，SM 可以获取第二个 Warp 并执行它。可以在不同的 Warp 之间进行调度，以确保最小化 SM 内发生的停顿次数。这就是为什么当块大小是 64 时，总执行时间比块大小是 32 时要短。不过将在后续章节中进行更多的性能分析。
